\chapter{Strings}\label{ch:strings}

String problems look diverse, but most of them reduce to a few questions. Does a pattern occur, and where? Are two substrings equal?
What is the longest common prefix of two suffixes? Which substrings are palindromes, and how many distinct substrings are there? This
chapter builds the standard tools in order of power: hashing, borders and the Z-function, tries and Aho--Corasick, the palindrome
tools, Lyndon words, and the two structures that capture all substrings at once, the suffix array and the suffix automaton.

\begin{chapterintro}{\gold\ \plat\ \ioi}{Chapter~\ref{ch:cpp} (strings in C++) and Chapter~\ref{ch:segtree} (Fenwick trees).}
  \item compare substrings in $\Oh{1}$ with hashing, and keep the hashes under updates;
  \item compute borders, periods and pattern matches with the prefix function and the Z-function, and turn KMP into an automaton;
  \item match many patterns at once with tries and Aho--Corasick;
  \item find palindromes with Manacher's algorithm and count them with the palindromic tree;
  \item factorise strings into Lyndon words and find minimal rotations;
  \item answer questions about all substrings with suffix arrays and suffix automata.
\end{chapterintro}

% ======================================================================================================================
\section[String hashing]{String hashing \hfill \gold}\label{sec:strhash}

A \defn{polynomial hash} maps a string to a number: $h(s) = s_0 B^{n-1} + s_1 B^{n-2} + \dots + s_{n-1} \bmod M$ for a base $B$ and a
modulus $M$. The prefix hashes $H_i = h(s_0 \dots s_{i-1})$ satisfy $H_{i+1} = H_i B + s_i$, and they give the hash of any substring in
$\Oh{1}$:
\[
  h\bigl(s[l, r)\bigr) = H_r - H_l \cdot B^{r - l} \pmod M ,
\]
because $H_r$ contains $H_l$ shifted by $r - l$ positions (Figure~\ref{fig:strhash}). Equal substrings have equal hashes. For a random
base, two \emph{different} strings of length at most $n$ collide with probability at most $n / M$, since their difference is a non-zero
polynomial in $B$ of degree below $n$, which has fewer than $n$ roots.

\begin{figure}[htbp]
\centering
\begin{tikzpicture}
  \node[blbl, anchor=east] at (-0.2, 0) {$s$};
  \arrs{H}{0}{0}{a/cool, b/cool, r/cool, a/cool, c/n, a/n, d/n, a/cool, b/cool, r/cool, a/cool}
  \idx{H}{11}
  \spanabove{H-1}{H-4}{$H_4$}
  \spanabove{H-8}{H-11}{$H_{11} - H_7 \cdot B^4$}
  \node[notel, anchor=west, text width=5.4cm] at (8.2, -0.1) {Both windows hash to the same number, so ``abra'' occurs at 0 and at 7:
    checked in $\Oh{1}$ after $\Oh{n}$ preprocessing, whatever the window length.};
\end{tikzpicture}
\caption{Substring hashes from prefix hashes: $h(s[7, 11)) = H_{11} - H_7 B^4$.}
\label{fig:strhash}
\end{figure}

\usetemplate{String hashing}{Ctrl+Alt+T, H}{cpstrhash}

\begin{pitfall}[Weak hashes get hacked]
With $M = 2^{64}$ (letting \cc{unsigned long long} overflow) the Thue--Morse strings collide for \emph{every} base, and with a fixed
small modulus and a fixed base, anybody can construct collisions in advance. The library uses $M = 2^{61} - 1$, a prime for which
multiplication stays fast, and a base drawn at run time: about $5 \cdot 10^{-13}$ per comparison of strings of length $10^6$, and
nothing to precompute against.
\end{pitfall}

\paragraph{Hashes under updates.} A hash is a sum of $s_i B^{\text{position}}$ terms, so a Fenwick tree over these terms keeps substring
hashes available while characters change: a point update replaces one term, and a range sum gives the hash, multiplied by a power of $B$.

\begin{problem}{Palindrome Queries (CSES)}
A string of length $n \le 2 \cdot 10^5$ receives $m \le 2 \cdot 10^5$ operations: change one character, or ask whether $s[a..b]$ is a
palindrome.
\solution Keep two Fenwick trees: one over $s_i B^i$ (the string read forwards) and one over $s_i B^{n-1-i}$ (backwards). The range
sums $F$ and $R$ over $[a, b]$ hash the substring and its reverse with different alignments. If the substring is a palindrome, then
$R = B^{n-1-a-b} F$, so the test is $F \cdot B^{n-1-b} = R \cdot B^{a}$. Each operation is $\Oh{\log n}$: 0.16\,s in total.
\end{problem}

\begin{code}
void solve() {  // StrHash (cpstrhash) is pasted above solve(): mul, base, MODH
    int n, m;
    string s;
    cin >> n >> m >> s;
    const uint64_t M = StrHash::MODH;
    vector<uint64_t> pw(n + 1, 1);
    rep(i, n) pw[i + 1] = StrHash::mul(pw[i], StrHash::base());
    // Fenwick trees over s[i] B^i (forward) and s[i] B^(n-1-i) (backward)
    vector<uint64_t> fw(n + 1), bw(n + 1);
    auto add = [&](vector<uint64_t>& f, int i, uint64_t v) {
        for (i++; i <= n; i += i & -i) f[i] = (f[i] + v) % M;
    };
    auto sum = [&](vector<uint64_t>& f, int i) {  // the first i entries
        uint64_t r = 0;
        for (; i > 0; i -= i & -i) r = (r + f[i]) % M;
        return r;
    };
    auto put = [&](int i, char c, bool remove) {  // add or remove letter c at i
        uint64_t a = StrHash::mul(c, pw[i]);
        uint64_t b = StrHash::mul(c, pw[n - 1 - i]);
        if (remove) a = M - a, b = M - b;
        add(fw, i, a), add(bw, i, b);
    };
    rep(i, n) put(i, s[i], false);
    while (m--) {
        int t;
        cin >> t;
        if (t == 1) {
            int k;
            char x;
            cin >> k >> x;
            put(k - 1, s[k - 1], true), s[k - 1] = x, put(k - 1, x, false);
        } else {
            int a, b;
            cin >> a >> b;
            a--, b--;
            uint64_t f = (sum(fw, b + 1) + M - sum(fw, a)) % M;
            uint64_t r = (sum(bw, b + 1) + M - sum(bw, a)) % M;
            // a palindrome reads the same both ways; align the powers
            yn(StrHash::mul(f, pw[n - 1 - b]) == StrHash::mul(r, pw[a]));
        }
    }
}
\end{code}

\begin{trick}[Hash plus binary search]
The longest common prefix of two suffixes is found by binary search on its length, comparing hashes: $\Oh{\log n}$. That gives
$\Oh{\log n}$ lexicographic comparison of any two substrings (compare the first character after the common prefix), the longest
palindrome centred at every position, the longest substring that occurs twice (binary search the length, with all hashes of that length
in a hash set), and suffix sorting in $\Oh{n \log^2 n}$ with an ordinary \cc{sort}.
\end{trick}

% ======================================================================================================================
\section[Borders and the prefix function]{Borders and the prefix function \hfill \plat}\label{sec:kmp}

A \defn{border} of a string is a proper prefix that is also a suffix: ``aba'' is a border of ``ababa''. The \defn{prefix function} $\pi[i]$
is the length of the longest border of $s[0..i]$ (Figure~\ref{fig:prefixfunc}). Two facts make it computable in $\Oh{n}$. First, a border
of $s[0..i]$, minus its last character, is a border of $s[0..i-1]$, so $\pi[i] \le \pi[i-1] + 1$. Second, the borders of $s[0..i-1]$ are
exactly $\pi[i-1]$, then $\pi[\pi[i-1] - 1]$, and so on, because a border of a border is a border. So try to extend the longest border
by $s[i]$; if that fails, fall back to the next shorter border and try again.

\begin{figure}[htbp]
\centering
\begin{tikzpicture}
  \node[blbl, anchor=east] at (-0.2, 0) {$s$};
  \arrs{P}{0}{0}{a/cool, b/cool, a/cool, b/n, a/n, c/n, a/warn, b/warn, a/warn, b/warn}
  \idx{P}{10}
  \node[blbl, anchor=east] at (-0.2, -1.2) {$\pi$};
  \arrs{Q}{0}{-1.2}{0/n, 0/n, 1/n, 2/n, 3/n, 0/n, 1/n, 2/n, 3/n, 4/hot}
  \spanabove{P-1}{P-4}{border ``abab''}
  \spanabove{P-7}{P-10}{suffix ``abab''}
  \node[notel, anchor=west, text width=5.6cm] at (7.5, -0.6) {$\pi[9] = 4$: the prefix ``abab'' (blue) equals the suffix (orange).
    The next border is $\pi[3] = 2$ (``ab''), then $\pi[1] = 0$: the borders of the whole string are 4 and 2.};
\end{tikzpicture}
\caption{The prefix function of ``ababacabab''.}
\label{fig:prefixfunc}
\end{figure}

The loop is linear by an amortised argument: the current border length $j$ grows by at most 1 per character, and every fallback
decreases it, so there are at most $n$ fallbacks in total. To find a pattern $p$ in a text $t$, compute $\pi$ of $p + \# + t$ with a
separator that occurs in neither: every position where $\pi = |p|$ ends an occurrence. This is the \textbf{Knuth--Morris--Pratt}
algorithm, $\Oh{|p| + |t|}$; String Matching (page~\pageref{wp:kmp}) writes it out in full.

\usetemplate{KMP (prefix function + pattern search)}{Ctrl+Alt+T, Shift+K}{cpkmp}

\begin{problem}{Finding Borders (CSES)}
Print the lengths of all borders of a string of length $n \le 10^6$, in increasing order.
\solution The longest border is $\pi[n-1]$, and following the chain $b \mapsto \pi[b - 1]$ visits all the others. 0.15\,s even for
``aaa\dots a'', which has $n - 1$ borders.
\end{problem}

\begin{code}
void solve() {  // prefix_function (cpkmp) is pasted above solve()
    string s;
    cin >> s;
    vi pi = prefix_function(s), borders;
    for (int b = pi.back(); b > 0; b = pi[b - 1]) borders.push_back(b);
    reverse(all(borders));  // the borders of a border are borders too
    print(borders);
}
\end{code}

\paragraph{The KMP automaton.} Precompute, for every matched length $j < |p|$ and every letter $c$, the matched length after reading $c$:
$j + 1$ if $p_j = c$, and otherwise the same transition from the next shorter border, $\mathit{go}(\pi[j-1], c)$. The result is a
deterministic automaton that tracks ``the longest suffix of the text read so far that is a prefix of $p$'' (Figure~\ref{fig:kmpauto}),
and DPs over strings can run on its states.

\begin{figure}[htbp]
\centering
\begin{tikzpicture}[x=2.1cm]
  \tikzset{ks/.style={vtx, minimum size=8mm}, kl/.style={font=\sffamily\scriptsize\bfseries, inner sep=1.5pt, fill=white}}
  \foreach \ii in {0, ..., 3} \node[ks] (K\ii) at (\ii, 0) {\ii};
  \node[ks, good, double, double distance=1pt] (K4) at (4, 0) {4};
  \foreach \ii/\jj/\klt in {0/1/A, 1/2/B, 2/3/A, 3/4/B} \draw[ded, line width=1pt, draw=fgblue] (K\ii) -- node[kl, text=fgblue] {\klt} (K\jj);
  \draw[ded] (K0) to[loop above, looseness=6] node[kl] {B} (K0);
  \draw[ded] (K1) to[loop above, looseness=6] node[kl] {A} (K1);
  \draw[ded] (K2) to[bend left=40] node[kl] {B} (K0);
  \draw[ded] (K3) to[bend left=40] node[kl] {A} (K1);
  \draw[ded] (K4) to[loop above, looseness=6] node[kl] {A, B} (K4);
  \node[notel, anchor=north, text width=11cm] at (2, -1.5) {The KMP automaton of ``ABAB'' over $\{A, B\}$. State $j$ means that the
    last $j$ letters read equal the first $j$ letters of the pattern. A mismatch falls back along the borders: after ``ABA'' a B
    completes the pattern, an A leaves ``A'' (state 1). For ``does the text contain $p$'' the final state is made absorbing.};
\end{tikzpicture}
\caption{A pattern-matching automaton built from the prefix function.}
\label{fig:kmpauto}
\end{figure}

\begin{problem}{Required Substring (CSES)}
Count the strings of length $n \le 1000$ over the alphabet A--Z that contain a given pattern of length $m \le 100$, modulo $10^9 + 7$.
\solution Count the complement: strings that never reach state $m$ of the KMP automaton. The DP over (length, state) has $nm$ states with
26 transitions each: 0.08\,s. The answer is $26^n$ minus the strings that avoid the pattern.
\end{problem}

\begin{code}
void solve() {  // prefix_function (cpkmp) is pasted above solve()
    int n;
    string p;
    cin >> n >> p;
    int m = sz(p);
    vi pi = prefix_function(p);
    vvi go(m, vi(26));  // the KMP automaton: matched length after a letter
    rep(j, m) rep(c, 26) {
        if (p[j] - 'A' == c) go[j][c] = j + 1;
        else go[j][c] = j ? go[pi[j - 1]][c] : 0;
    }
    vi dp(m), nd(m);  // dp[j]: strings so far whose longest match with p is j
    dp[0] = 1;
    ll total = 1;
    rep(i, n) {
        fill(all(nd), 0);
        rep(j, m) rep(c, 26) {
            int to = go[j][c];
            if (to < m) nd[to] = (nd[to] + dp[j]) % MOD;
        }
        swap(dp, nd);  // strings that complete p are dropped: they contain it
        total = total * 26 % MOD;
    }
    ll avoid = 0;
    for (ll v : dp) avoid = (avoid + v) % MOD;
    print((total - avoid + MOD) % MOD);
}
\end{code}

% ======================================================================================================================
\section[The Z-function]{The Z-function \hfill \plat}\label{sec:zfunc}

The \defn{Z-function} $z[i]$ is the length of the longest common prefix of $s$ and its suffix starting at $i$. The algorithm keeps the
rightmost \emph{Z-box} $[l, r)$ found so far, a window that matches the prefix $s[0, r - l)$ (Figure~\ref{fig:zfunc}). For a position $i$
inside it, the window says that $s[i, r)$ equals $s[i - l, r - l)$, so $z[i] \ge \min(r - i, z[i - l])$ for free; only characters
beyond $r$ are compared explicitly, and each such comparison moves $r$ to the right. Hence $\Oh{n}$ in total.

\begin{figure}[htbp]
\centering
\begin{tikzpicture}
  \node[blbl, anchor=east] at (-0.2, 0) {$s$};
  \arrs{Z}{0}{0}{a/cool, a/cool, b/cool, x/cool, a/warn, a/warn, b/warn, x/warn, c/n, a/n, a/n, b/n}
  \idx{Z}{12}
  \node[blbl, anchor=east] at (-0.2, -1.2) {$z$};
  \arrs{Y}{0}{-1.2}{0/n, 1/n, 0/n, 0/n, 4/hot, 1/n, 0/n, 0/n, 0/n, 3/n, 1/n, 0/n}
  \spanabove{Z-5}{Z-8}{Z-box $[4, 8)$}
  \node[notel, anchor=west, text width=4.6cm] at (8.9, -0.6) {$z[4] = 4$: ``aabx'' at 4 matches the prefix. Inside the box, $z[5]$ starts
    from $z[1] = 1$ without comparing anything.};
\end{tikzpicture}
\caption{The Z-function of ``aabxaabxcaab''.}
\label{fig:zfunc}
\end{figure}

\usetemplate{Z-function}{Ctrl+Alt+T, Z}{cpz}

The Z-function and the prefix function carry the same information and convert into each other in $\Oh{n}$; use whichever reads more
naturally. Matching a pattern is $z$ of $p + \# + t$; periods read off $z$ directly.

\begin{problem}{Finding Periods (CSES)}
Print all periods of a string of length $n \le 10^6$: the lengths $p$ such that $s_i = s_{i+p}$ for every valid $i$.
\solution $p < n$ is a period exactly when the suffix starting at $p$ is a prefix of $s$, that is $z[p] = n - p$; $n$ itself always is.
0.11\,s.
\end{problem}

\begin{code}
void solve() {  // z_function (cpz) is pasted above solve()
    string s;
    cin >> s;
    int n = sz(s);
    vi z = z_function(s), periods;
    rep(p, 1, n) if (z[p] == n - p) periods.push_back(p);  // s[p..) repeats s
    periods.push_back(n);
    print(periods);
}
\end{code}

\begin{fact}[Periods and borders]
$p$ is a period of $s$ exactly when $s$ has a border of length $n - p$, so the smallest period is $n - \pi[n-1]$, and $s$ is a power of a
shorter string exactly when that period divides $n$. \textbf{Fine--Wilf}: if $p$ and $q$ are periods and $p + q - \gcd(p, q) \le n$, then
$\gcd(p, q)$ is a period too.
\end{fact}

% ======================================================================================================================
\section[Tries]{Tries \hfill \gold}\label{sec:trie}

A \defn{trie} stores a set of strings as a rooted tree whose edges are letters: every string is the path from the root to a node, and
strings with a common prefix share that part of the path (Figure~\ref{fig:trie}). Insertion and lookup cost $\Oh{|s|}$, and a node can
store counts, such as the number of words that pass through it (``how many words start with $p$''). The binary trie of
Section~\ref{sec:xor} is the same structure over the bits of numbers.

\begin{figure}[htbp]
\centering
\begin{tikzpicture}[level distance=11mm, level 1/.style={sibling distance=26mm}, level 2/.style={sibling distance=12mm}]
  \tikzset{tw/.style={vtxs}, te/.style={font=\ttfamily\small, fill=white, inner sep=1pt}}
  \node[tw] (R) {}
    child {node[tw, hot] (A) {}
      child {node[tw, hot] (AB) {}
        child {node[tw] (ABA) {}
          child {node[tw, hot] (ABAB) {} edge from parent[ed] node[te, right] {b}}
          edge from parent[ed] node[te, right] {a}}
        edge from parent[ed] node[te, right] {b}}
      edge from parent[ed] node[te, left] {a}}
    child {node[tw, hot] (B) {} edge from parent[ed] node[te, right] {b}};
  \node[notel, anchor=west, text width=6.6cm] at (2.4, -1.9) {The trie of the words a, ab, abab and b (red: a word ends). Splitting
    $s = $ abab: from position 0 the walk passes the ends of a, ab and abab, so $\mathit{dp}[0] = \mathit{dp}[1] + \mathit{dp}[2] +
    \mathit{dp}[4] = 2 + 2 + 1 = 5$ ways.};
\end{tikzpicture}
\caption{A trie, and the walks of Word Combinations.}
\label{fig:trie}
\end{figure}

\usetemplate{Trie / binary trie}{Ctrl+Alt+T, Shift+T / \textasciigrave}{cptrie, cpxortrie}

\begin{problem}{Word Combinations (CSES)}
In how many ways can a string of length $n \le 5000$ be split into words from a dictionary of $k \le 10^5$ distinct words (total length
up to $10^6$)? Answer modulo $10^9 + 7$.
\solution $\mathit{dp}[i]$, the number of ways to split $s[i..)$, is the sum of $\mathit{dp}[j + 1]$ over the dictionary words $s[i..j]$.
Walking the trie along $s[i..]$ finds all of them in one pass, stopping as soon as no word continues: $\Oh{n \cdot L}$ for the longest
word length $L$. 0.15\,s.
\end{problem}

\begin{code}
void solve() {  // Trie (cptrie) is pasted above solve()
    string s;
    int k;
    cin >> s >> k;
    Trie t;
    rep(i, k) {
        string w;
        cin >> w;
        t.insert(w);
    }
    int n = sz(s);
    vi dp(n + 1);  // dp[i]: the ways to split s[i..)
    dp[n] = 1;
    per(i, n) {
        int u = 0;  // walk the trie along s[i..]; word ends split
        for (int j = i; j < n && (u = t.nxt[u][s[j] - 'a']); j++)
            if (t.ends[u]) dp[i] = (dp[i] + dp[j + 1]) % MOD;
    }
    print(dp[0]);
}
\end{code}

% ======================================================================================================================
\section[Aho--Corasick]{Aho--Corasick \hfill \ioi}\label{sec:aho}

To find many patterns in a text at once, build their trie and add \defn{suffix links}: the link of a node points to the node of the
longest proper suffix of its string that is also in the trie (Figure~\ref{fig:aho}). This generalises the prefix function from one
pattern to a trie, and the links are computed the same way, in BFS order: the link of $u$'s child by letter $c$ is the child by $c$ of
$u$'s link, falling back further when it does not exist. Filling in every missing child with that fallback turns the trie into a
complete automaton, and the text is then read in one pass, one transition per character.

\begin{figure}[htbp]
\centering
\begin{tikzpicture}
  \tikzset{an/.style={rectangle, rounded corners=2pt, draw=fgink!80, fill=fgfill, line width=0.7pt, font=\ttfamily\scriptsize,
    minimum width=8mm, minimum height=5.6mm, inner sep=1.5pt}, ao/.style={an, good},
    al/.style={->, line width=0.8pt, draw=fgverm, dashed, shorten >=1pt}}
  \node[an] (r) at (0, 0) {$\varepsilon$};
  \node[an] (h) at (-2.6, -1.1) {h};
  \node[an] (s) at (2.6, -1.1) {s};
  \node[ao] (he) at (-3.6, -2.2) {he};
  \node[an] (hi) at (-1.6, -2.2) {hi};
  \node[an] (her) at (-3.6, -3.3) {her};
  \node[ao] (hers) at (-3.6, -4.4) {hers};
  \node[ao] (his) at (-1.6, -3.3) {his};
  \node[an] (sh) at (2.6, -2.2) {sh};
  \node[ao] (she) at (2.6, -3.3) {she};
  \draw[ed] (r) -- (h) (r) -- (s) (h) -- (he) (h) -- (hi) (he) -- (her) (her) -- (hers) (hi) -- (his) (s) -- (sh) (sh) -- (she);
  \draw[al] (sh) to[bend right=12] (h);
  \draw[al] (she) to[bend left=6] (he);
  \draw[al] (his) to[bend right=15] (s);
  \draw[al] (hers.east) to[bend right=30] (s.south);
  \node[notel, anchor=west, text width=4.4cm] at (4.2, -2.3) {The trie of he, she, his and hers (green: a pattern ends). Red dashed
    arrows are the suffix links that do not lead to the root. Reading ``she'' ends in node she, whose link reaches he: both patterns
    occur there.};
\end{tikzpicture}
\caption{The Aho--Corasick automaton: a trie plus suffix links.}
\label{fig:aho}
\end{figure}

A position of the text can end several patterns: those of the current node and of every node on its suffix-link path. Rather than walk
that path at every character, count how often each node is the current one, and afterwards push the counts up the suffix links, in
reverse BFS order: then every node's count is the number of occurrences of its string. The whole algorithm costs $\Oh{|\Sigma| \cdot
\text{total pattern length} + |t|}$.

\usetemplate{Aho--Corasick}{Ctrl+Alt+Y, Shift+A}{cpaho}

\begin{problem}{Counting Patterns (CSES)}
For a text of length $10^5$ and $k$ patterns of total length $5 \cdot 10^5$, print how many times each pattern occurs.
\solution Aho--Corasick with counts pushed up the suffix links; identical patterns share a node. 0.08\,s for $10^5$ patterns.
\end{problem}

\begin{code}
void solve() {  // AhoCorasick (cpaho) is pasted above solve()
    string text;
    int k;
    cin >> text >> k;
    AhoCorasick ac;
    vi id(k);
    rep(i, k) {
        string p;
        cin >> p;
        id[i] = ac.add(p);  // equal patterns share a node
    }
    ac.build();
    vector<ll> occ = ac.occurrences(text);  // counts pushed up the suffix links
    rep(i, k) print(occ[id[i]]);
}
\end{code}

\begin{recognize}
Many patterns in one text (``count the dictionary words in a document''), and DPs over strings that must avoid or must contain a set of
words: the automaton's states become the DP states, and ``strings of length $n$ with no forbidden word'' are walks of length $n$ that
never enter a terminal state, a matrix power when $n$ is huge.
\end{recognize}

% ======================================================================================================================
\section[Palindromes]{Palindromes \hfill \plat}\label{sec:palindromes}

\paragraph{Manacher's algorithm.} For every centre $i$, let $d_1[i]$ be the number of odd palindromes centred at $i$ (the radius of the
longest one, counting the centre). Manacher's algorithm computes all of them in $\Oh{n}$ with the idea of the Z-box: keep the palindrome
$[l, r]$ that reaches furthest to the right. A centre $i$ inside it has a mirror image $j = l + r - i$, whose radius is already known, and
the reflection guarantees $d_1[i] \ge \min(d_1[j], r - i + 1)$ (Figure~\ref{fig:manacher}). Only characters beyond $r$ are compared, and
each comparison moves $r$. Even palindromes get a second array, or the string is interleaved with separators.

\begin{figure}[htbp]
\centering
\begin{tikzpicture}
  \arrs{M}{0}{0}{a/n, b/n, a/n, c/hot, a/n, b/n, a/n, d/n, a/n, b/n}
  \idx{M}{10}
  \draw[line width=1.2pt, draw=fgblue] (M-1.north) .. controls ($(M-1.north)+(0, 1.3)$) and ($(M-7.north)+(0, 1.3)$) .. (M-7.north);
  \node[font=\sffamily\scriptsize, text=fgblue] at ($(M-4.north)+(0, 1.1)$) {``abacaba'' centred at 3};
  \draw[line width=1pt, draw=fggreen] ($(M-1.south)+(0, -0.42)$) .. controls ($(M-1.south)+(0, -1.05)$) and ($(M-3.south)+(0, -1.05)$) .. ($(M-3.south)+(0, -0.42)$);
  \draw[line width=1pt, draw=fggreen] ($(M-5.south)+(0, -0.42)$) .. controls ($(M-5.south)+(0, -1.05)$) and ($(M-7.south)+(0, -1.05)$) .. ($(M-7.south)+(0, -0.42)$);
  \draw[->, line width=0.8pt, draw=fgink!60] ($(M-2.south)+(0, -1.05)$) to[bend right=20] node[below, lbl] {mirror} ($(M-6.south)+(0, -1.05)$);
  \node[notel, anchor=west, text width=5.0cm] at (7.4, 0) {Inside the big palindrome, the centre 5 mirrors the centre 1, so it
    inherits the palindrome ``aba'' (green) without a comparison. Here the radii are $d_1 = 1, 2, 1, 4, 1, 2, 1, 3, 1, 1$.};
\end{tikzpicture}
\caption{Manacher's algorithm reuses the mirror image inside the rightmost palindrome.}
\label{fig:manacher}
\end{figure}

\usetemplate{Manacher}{Ctrl+Alt+Y, 6}{cpmanacher}

\begin{problem}{Longest Palindrome (CSES)}
Find a longest palindromic substring of a string of length $n \le 10^6$.
\solution The largest radius among both arrays: 0.08\,s. The total number of palindromic substrings is $\sum_i d_1[i] + \sum_i d_2[i]$.
\end{problem}

\begin{code}
void solve() {  // manacher (cpmanacher) is pasted above solve()
    string s;
    cin >> s;
    auto [start, len] = longest_palindrome(s);
    print(s.substr(start, len));
}
\end{code}

\paragraph{The palindromic tree.} A string of length $n$ has at most $n$ \emph{distinct} palindromic substrings, because each new character
creates at most one new one: the longest palindromic suffix of the new prefix (a shorter palindromic suffix is a mirror image inside the
longest one, so it occurred earlier). The \defn{palindromic tree} (eertree) has one node per distinct palindrome. Its edges add a letter on
both sides ($X \to cXc$), and two roots start the odd and the even palindromes: an imaginary one of length $-1$ (adding $c$ on both
sides gives ``$c$'') and the empty palindrome of length 0. Every node also has a suffix link to its longest proper palindromic suffix
(Figure~\ref{fig:eertree}). Appending a character means walking the suffix links from the current longest palindromic suffix until one
can be extended by the new character, and the walks cost $\Oh{n}$ in total.

\begin{figure}[htbp]
\centering
\begin{tikzpicture}
  \tikzset{en/.style={rectangle, rounded corners=2pt, draw=fgink!80, fill=fgfill, line width=0.7pt, font=\ttfamily\small,
    minimum width=10mm, minimum height=6mm}, el/.style={->, line width=0.8pt, draw=fgverm, dashed, shorten >=1pt},
    et/.style={font=\ttfamily\scriptsize\bfseries, fill=white, inner sep=1pt}}
  \node[en, dim] (m1) at (0, 0) {$-1$};
  \node[en, dim] (z0) at (4.2, 0) {$\varepsilon$};
  \node[en] (pa) at (-1, -1.3) {a};
  \node[en] (pb) at (1, -1.3) {b};
  \node[en] (pbb) at (4.2, -1.3) {bb};
  \node[en, good] (pabba) at (4.2, -2.6) {abba};
  \draw[ded] (m1) -- node[et, left] {a} (pa);
  \draw[ded] (m1) -- node[et, right] {b} (pb);
  \draw[ded] (z0) -- node[et, right] {b} (pbb);
  \draw[ded] (pbb) -- node[et, right] {a} (pabba);
  \draw[el] (pa.south) to[bend right=35] (z0.south west);
  \draw[el] (pb) to[bend left=10] (z0);
  \draw[el] (pbb) -- (pb);
  \draw[el] (pabba) to[bend left=25] (pa.south);
  \draw[el] (z0) -- (m1);
  \node[notel, anchor=west, text width=5.2cm] at (5.7, -1.3) {The palindromic tree of ``abba'': four distinct palindromes. Solid edges
    put a letter on both sides; red dashed edges are suffix links (abba $\to$ a, bb $\to$ b, single letters $\to \varepsilon$, and
    $\varepsilon \to -1$).};
\end{tikzpicture}
\caption{The palindromic tree has one node per distinct palindrome.}
\label{fig:eertree}
\end{figure}

\usetemplate{Palindromic tree (eertree)}{Ctrl+Alt+U, Shift+E}{cpeertree}

\begin{problem}{Palindromes (APIO 2014)}
For a string of length $n \le 3 \cdot 10^5$, find the largest value of (length $\times$ number of occurrences) over its palindromic
substrings.
\solution Count, for every node, how often it was the longest palindromic suffix of a prefix; a palindrome also occurs wherever a longer
palindrome ending at the same place has it as a suffix, so push the counts down the suffix links, from longer to shorter. 0.07\,s.
\end{problem}

\begin{code}
void solve() {  // Eertree (cpeertree) is pasted above solve()
    string s;
    cin >> s;
    Eertree t;
    for (char c : s) t.add(c);
    vector<ll> occ = t.occurrences();  // per distinct palindrome
    ll best = 0;
    rep(v, 2, sz(t.t)) chmax(best, t.t[v].len * occ[v]);  // 0, 1 are the roots
    print(best);
}
\end{code}

% ======================================================================================================================
\section[Lyndon words and the minimal rotation]{Lyndon words and the minimal rotation \hfill \ioi}\label{sec:lyndon}

A \defn{Lyndon word} is a string that is strictly smaller than all of its proper rotations, such as ``aab'' or ``abc'' (but not ``aba'' or
``abab''). The \textbf{Chen--Fox--Lyndon theorem}: every string has a unique factorisation $s = w_1 w_2 \cdots w_k$ into Lyndon words with
$w_1 \ge w_2 \ge \dots \ge w_k$ (Figure~\ref{fig:lyndon}). \textbf{Duval's algorithm} computes it in $\Oh{n}$ with three pointers. It keeps a
prefix of the remaining string of the form $w w \cdots w w'$, where $w$ is a Lyndon word and $w'$ a proper prefix of $w$, and compares the
next character $s_j$ with the character $s_k$ one period earlier. If they are equal, the periodic pattern continues; if $s_j > s_k$, the
whole prefix together with $s_j$ becomes a single Lyndon word; if $s_j < s_k$, the copies of $w$ are output as factors and the scan
restarts after them.

\begin{figure}[htbp]
\centering
\begin{tikzpicture}
  \tikzset{lf/.style={rectangle, rounded corners=2pt, draw=fgblue, fill=fgsky!25, line width=0.8pt, font=\ttfamily\small,
    minimum height=7mm, inner sep=3pt}}
  \node[lf] (f1) at (0, 0) {b};
  \node[lf, anchor=west] (f2) at ($(f1.east)+(0.25, 0)$) {ac};
  \node[lf, anchor=west] (f3) at ($(f2.east)+(0.25, 0)$) {abc};
  \node[lf, anchor=west] (f4) at ($(f3.east)+(0.25, 0)$) {ab};
  \node[lf, anchor=west] (f5) at ($(f4.east)+(0.25, 0)$) {a};
  \foreach \fa/\fb in {f1/f2, f2/f3, f3/f4, f4/f5}
    \node[font=\small, text=fgverm] at ($(\fa.east)!0.5!(\fb.west)$) {$\ge$};
  \node[notel, anchor=west, text width=6.6cm] at ($(f5.east)+(0.6, 0)$) {The Lyndon factorisation of ``bacabcaba'': five Lyndon words in
    non-increasing order. The minimal rotation of a string starts at the last factor of $s + s$ that begins in the first half: for
    ``abbaabbab'' it is ``aabbababb''.};
\end{tikzpicture}
\caption{The unique factorisation into non-increasing Lyndon words.}
\label{fig:lyndon}
\end{figure}

The \textbf{minimal rotation} of $s$, a canonical form for cyclic strings and necklaces, starts at the beginning of the last Lyndon factor
of $s + s$ that starts in the first half, so Duval's algorithm finds it in $\Oh{n}$ too.

\usetemplate{Lyndon factorization (Duval) and minimal rotation}{Ctrl+Alt+U, Y}{cplyndon}

\begin{problem}{Minimal Rotation (CSES)}
Print the lexicographically smallest rotation of a string of length $n \le 10^6$.
\solution Duval on $s + s$: 0.08\,s.
\end{problem}

\begin{code}
void solve() {  // min_rotation (cplyndon) is pasted above solve()
    string s;
    cin >> s;
    int k = min_rotation(s);  // Duval's algorithm on s + s
    print(s.substr(k) + s.substr(0, k));
}
\end{code}

% ======================================================================================================================
\section[Suffix arrays]{Suffix arrays \hfill \ioi}\label{sec:suffixarray}

The \defn{suffix array} lists the starting positions of all suffixes in lexicographic order, and the \defn{LCP array} stores the longest
common prefix of every two neighbours in that order (Figure~\ref{fig:suffixarray}). Every substring is a prefix of some suffix, and the
suffixes that start with a given string form one contiguous block of the order, so the two arrays answer most questions about all
substrings at once.

\begin{figure}[htbp]
\centering
\begin{tikzpicture}[y=0.62cm]
  \foreach \rk/\sav/\suf/\lc in {0/5/a/0, 1/3/ana/1, 2/1/anana/3, 3/0/banana/0, 4/4/na/0, 5/2/nana/2} {
    \node[lbl] at (0, -\rk) {\rk};
    \node[font=\ttfamily\small] at (0.9, -\rk) {\sav};
    \node[font=\ttfamily\small] at (1.8, -\rk) {\lc};
    \fill[fgsky!40] (2.6, -\rk - 0.3) rectangle ++(0.183*\lc, 0.6);
    \node[font=\ttfamily\small, anchor=west, inner sep=0pt] at (2.6, -\rk) {\suf};
  }
  \node[blbl] at (0, 0.9) {rank};
  \node[blbl] at (0.9, 0.9) {sa};
  \node[blbl] at (1.8, 0.9) {lcp};
  \node[blbl, anchor=west] at (2.6, 0.9) {suffix};
  \node[notel, anchor=west, text width=6.8cm] at (5.2, -2.5) {The sorted suffixes of ``banana''. The shaded part of each suffix is its
    common prefix with the previous one (the lcp column, here aligned with the later suffix). The distinct substrings are the
    prefixes that are not shaded: $21 - (1 + 3 + 0 + 0 + 2) = 15$.};
\end{tikzpicture}
\caption{The suffix array and the LCP array of ``banana''.}
\label{fig:suffixarray}
\end{figure}

\paragraph{Construction.} Sort the cyclic shifts of $s$ with a sentinel appended by their first $2^k$ characters, for $k = 0, 1, 2,
\dots$ After round $k$ every shift has a rank among the strings of length $2^k$. The first $2^{k+1}$ characters of a shift are two
halves of length $2^k$ whose ranks are already known, so round $k + 1$ sorts pairs of ranks, a radix sort in $\Oh{n}$. After $\log n$
rounds all suffixes are ordered: $\Oh{n \log n}$. The LCP array follows in $\Oh{n}$ by \textbf{Kasai's algorithm}: process the suffixes
in order of their starting position; when the suffix at $i$ has a common prefix of length $h$ with its successor in the order, the suffix
at $i + 1$ has one of length at least $h - 1$ with its own successor, so the comparison resumes from $h - 1$.

\usetemplate{Suffix array + LCP}{Ctrl+Alt+Y, 4}{cpsa}

\begin{problem}{Distinct Substrings (CSES)}
Count the distinct substrings of a string of length $n \le 10^5$.
\solution Every substring is a prefix of a suffix, and walking the suffixes in sorted order, the prefixes of each suffix that were not
already prefixes of the previous one are exactly those longer than their common prefix. The answer is $\frac{n(n+1)}{2} - \sum
\mathit{lcp}$: 0.08\,s.
\end{problem}

\begin{code}
void solve() {  // suffix_array, lcp_array (cpsa) are pasted above solve()
    string s;
    cin >> s;
    ll n = sz(s), ans = n * (n + 1) / 2;
    vi sa = suffix_array(s), lcp = lcp_array(s, sa);
    for (ll x : lcp) ans -= x;  // prefixes shared with the previous suffix
    print(ans);
}
\end{code}

\begin{problem}{Substring Order I (CSES)}
Find the $k$-th smallest distinct substring of a string of length $n \le 10^5$.
\solution In sorted order, the suffix of rank $r$ contributes the new substrings of lengths $\mathit{lcp}[r-1] + 1, \dots, n - \mathit{sa}[r]$,
and they come in increasing order. Skip whole suffixes until the $k$-th one falls inside: 0.08\,s.
\end{problem}

\begin{code}
void solve() {  // suffix_array, lcp_array (cpsa) are pasted above solve()
    string s;
    ll k;
    cin >> s >> k;
    int n = sz(s);
    vi sa = suffix_array(s), lcp = lcp_array(s, sa);
    rep(r, n) {
        ll shared = r ? lcp[r - 1] : 0;  // counted already at rank r - 1
        ll fresh = n - sa[r] - shared;  // lengths shared + 1 .. n - sa[r]
        if (k <= fresh) {
            print(s.substr(sa[r], shared + k));
            return;
        }
        k -= fresh;
    }
}
\end{code}

\begin{problem}{Substring Distribution (CSES)}
For every length from 1 to $n$, count the distinct substrings of that length ($n \le 10^5$).
\solution Each rank contributes one new substring of every length in an interval; add the intervals with a difference array. 0.09\,s.
\end{problem}

\begin{code}
void solve() {  // suffix_array, lcp_array (cpsa) are pasted above solve()
    string s;
    cin >> s;
    int n = sz(s);
    vi sa = suffix_array(s), lcp = lcp_array(s, sa), diff(n + 2);
    rep(r, n) {  // rank r: one new substring per length shared + 1 .. n - sa[r]
        ll shared = r ? lcp[r - 1] : 0;
        diff[shared + 1]++, diff[n - sa[r] + 1]--;
    }
    vi cnt(n);
    ll run = 0;
    rep(len, 1, n + 1) cnt[len - 1] = run += diff[len];
    print(cnt);
}
\end{code}

\begin{trick}[More from suffix arrays]
\begin{itemize}
  \item The LCP of any two suffixes is the minimum of the LCP array between their ranks: a sparse table answers it in $\Oh{1}$.
  \item The longest repeated substring is the largest LCP value (Repeating Substring).
  \item The longest common substring of $s$ and $t$: build the array of $s + \# + t$ and take the largest LCP between neighbours that come
        from different strings.
  \item A pattern's occurrences form a block of the order, found by binary search in $\Oh{|p| \log n}$.
  \item The suffix array with its LCP array is a flattened \emph{suffix tree}, the compressed trie of all suffixes: its inner nodes are the
        intervals of the order whose minimum LCP is the node's depth, so suffix-tree algorithms run on it with a stack.
\end{itemize}
\end{trick}

% ======================================================================================================================
\section[Suffix automata]{Suffix automata \hfill \ioi}\label{sec:sam}

The \defn{suffix automaton} of $s$ is the smallest deterministic automaton that accepts exactly the substrings of $s$. Group the substrings
by their set of end positions in $s$ (their \emph{endpos}): one state per group. Within a group the strings are suffixes of each other with
consecutive lengths, so a state is described by its longest string, of length $\mathit{len}$, and by a \emph{suffix link} to the group
of the next shorter suffix (Figure~\ref{fig:sam}). The automaton has at most $2n - 1$ states and $3n - 4$ transitions, and it is built
online, one character at a time, in amortised $\Oh{1}$ per character: the new state is linked to existing ones along the suffix links,
and when a transition would put strings with different endpos sets into one state, that state is split by a \emph{clone}.

\begin{figure}[htbp]
\centering
\begin{tikzpicture}
  \tikzset{ss/.style={rectangle, rounded corners=3pt, draw=fgink!80, fill=fgfill, line width=0.7pt, font=\ttfamily\scriptsize,
    minimum width=12mm, minimum height=7mm, align=center}, sl/.style={->, line width=0.8pt, draw=fgverm, dashed, shorten >=1pt},
    st/.style={font=\ttfamily\scriptsize\bfseries, fill=white, inner sep=1pt}}
  \node[ss, dim] (S0) at (0, 0) {$\varepsilon$};
  \node[ss] (S1) at (2.4, 0) {a};
  \node[ss] (S2) at (4.8, 0) {ab};
  \node[ss, good] (S3) at (7.2, 0) {abb\\bb};
  \node[ss, warn] (S4) at (3.6, -1.8) {b};
  \draw[ded] (S0) -- node[st, above] {a} (S1);
  \draw[ded] (S1) -- node[st, above] {b} (S2);
  \draw[ded] (S2) -- node[st, above] {b} (S3);
  \draw[ded] (S0) -- node[st, below left] {b} (S4);
  \draw[ded] (S4) -- node[st, below right] {b} (S3);
  \draw[sl] (S1) to[bend right=40] (S0);
  \draw[sl] (S2) to[bend right=10] (S4);
  \draw[sl] (S3) to[bend left=25] (S4);
  \draw[sl] (S4) to[bend left=30] (S0);
  \node[notel, anchor=west, text width=4.4cm] at (8.2, -0.9) {The suffix automaton of ``abb''. State ``b'' (orange) is the clone
    created by the second b, when ``b'' started to end at two positions. Red dashed arrows are suffix links. Distinct substrings:
    $\sum (\mathit{len} - \mathit{len}_{\mathrm{link}}) = 1 + 1 + 2 + 1 = 5$.};
\end{tikzpicture}
\caption{A suffix automaton: one state per endpos class.}
\label{fig:sam}
\end{figure}

\usetemplate{Suffix automaton}{Ctrl+Alt+Y, 5}{cpsam}

\begin{fact}[What a suffix automaton answers]
The distinct substrings number $\sum_v (\mathit{len}(v) - \mathit{len}(\mathit{link}(v)))$. The number of occurrences of the strings of
a state is the size of its endpos set: give every prefix's own state a count of 1 and add the counts up the suffix links, from the
longest states down. The longest common substring of $s$ and $t$ is found by walking $t$ through the automaton of $s$, keeping the length
of the current match. Since transitions always lead to longer states, sorting the states by $\mathit{len}$ orders the automaton
topologically, and DPs over paths (the number of distinct substrings starting with a given prefix) run in that order.
\end{fact}

\begin{problem}{Substring Order II (CSES)}
Find the $k$-th smallest substring of a string of length $n \le 10^5$, counting every occurrence separately.
\solution Compute the occurrence count of every state, then $\mathit{total}(v)$, the number of substring occurrences on the paths that
start at $v$: its own occurrences plus the totals of its successors. Walk from the start, choosing at each step the smallest letter whose
total still reaches $k$: 0.12\,s.
\end{problem}

\begin{code}
void solve() {  // SuffixAutomaton (cpsam) is pasted above solve()
    string s;
    ll k;
    cin >> s >> k;
    SuffixAutomaton sam(s);
    auto& st = sam.st;
    int m = sz(st);
    vi occ(m), total(m), order(m);
    int u = 0;
    for (char ch : s) u = st[u].next[ch - 'a'], occ[u] = 1;  // prefix states
    iota(all(order), 0);
    sort(all(order), [&](int a, int b) { return st[a].len > st[b].len; });
    for (int v : order)
        if (v) occ[st[v].link] += occ[v];  // endpos sizes, longest states first
    for (int v : order) {  // substrings with repetitions on paths from v
        total[v] = v ? occ[v] : 0;
        for (int w : st[v].next)
            if (w != -1) total[v] += total[w];
    }
    string ans;
    for (int v = 0; k > 0;) {
        rep(c, 26) {
            int w = st[v].next[c];
            if (w == -1) continue;
            if (total[w] < k) {
                k -= total[w];
                continue;
            }
            ans += char('a' + c), k -= occ[w], v = w;  // itself comes first
            break;
        }
    }
    print(ans);
}
\end{code}

% ======================================================================================================================
\section[Choosing a string tool]{Choosing a string tool \hfill \gold}\label{sec:stringtools}

The tools overlap, and the simplest one that answers the question is usually the right choice. Figure~\ref{fig:stringflow} sorts them
by the question asked.

\begin{figure}[htbp]
\centering
\begin{tikzpicture}[node distance=2.6mm]
  \node[fq=5.2cm] (q1) {many patterns at once?};
  \node[fend=4.6cm, right=12mm of q1] (a1) {Aho--Corasick};
  \node[fq=5.2cm, below=of q1] (q2) {palindromes?};
  \node[fend=4.6cm, right=12mm of q2] (a2) {Manacher; palindromic tree for distinct ones};
  \node[fq=5.2cm, below=of q2] (q3) {equality of arbitrary substrings, or updates?};
  \node[fend=4.6cm, right=12mm of q3] (a3) {hashing (+ Fenwick tree for updates)};
  \node[fq=5.2cm, below=of q3] (q4) {one pattern, borders, periods?};
  \node[fend=4.6cm, right=12mm of q4] (a4) {prefix function or Z-function};
  \node[fq=5.2cm, below=of q4] (q5) {order, distinct substrings, LCPs?};
  \node[fend=4.6cm, right=12mm of q5] (a5) {suffix array + LCP};
  \node[fend=5.2cm, below=of q5] (a6) {suffix automaton (occurrences, online)};
  \foreach \qq/\qa in {q1/a1, q2/a2, q3/a3, q4/a4, q5/a5} \draw[farr] (\qq) -- node[fyes, above] {yes} (\qa);
  \foreach \qq/\nn in {q1/q2, q2/q3, q3/q4, q4/q5, q5/a6} \draw[farr] (\qq) -- node[fno, right] {no} (\nn);
\end{tikzpicture}
\caption{Which string tool to use.}
\label{fig:stringflow}
\end{figure}

\begin{summary}[Chapter summary]
\begin{itemize}
  \item Polynomial hashes modulo $2^{61} - 1$ with a random base compare substrings in $\Oh{1}$; Fenwick trees keep them under updates.
  \item The prefix function finds borders, periods and matches in $\Oh{n}$ by falling back along borders; it defines the KMP automaton
        for DPs over strings. The Z-function gives the same information from the other side.
  \item Tries share common prefixes; Aho--Corasick adds suffix links and matches many patterns in one pass.
  \item Manacher finds all palindromes around every centre in $\Oh{n}$; the palindromic tree has one node per distinct palindrome.
  \item Duval's algorithm factorises a string into non-increasing Lyndon words and finds the minimal rotation in $\Oh{n}$.
  \item Suffix arrays with LCP arrays count and order all substrings; suffix automata recognise them, one state per endpos class, and
        count occurrences.
\end{itemize}
\end{summary}

\begin{practice}
\ppgroup{String algorithms}
\begin{pplist}
  \item Word Combinations
  \item String Matching
  \item Finding Borders
  \item Finding Periods
  \item Minimal Rotation
  \item Longest Palindrome
  \item Required Substring
  \item Palindrome Queries
  \item Finding Patterns
  \item Counting Patterns
  \item Pattern Positions
  \item Distinct Substrings
  \item Repeating Substring
  \item String Functions
  \item Substring Order I
  \item Substring Order II
  \item Substring Distribution
\end{pplist}
\end{practice}
